\begin{center}
\hypertarget{native-042-article025}{}\pdfbookmark[2]{25.}{article-025}25.
\end{center}
Ad solutionem problematis inversi, e tempore anomaliam veram radiumque vectorem determinare, primo ex $N=\lambda kb^{-\frac{3}{2}}t$ per aequationem XI elicienda est quantitas auxiliaris $u$ vel $F$. Solutio huius aequationis transcendentis tentando perficienda erit, et per artificia iis quae in art. 11 exposuimus analoga abbreviari poterit. Haec autem fusius explicare supersedemus: neque enim operae pretium esse videtur, praecepta pro motu hyperbolico in coelis vix unquam fortasse se oblaturo aeque anxie expolire ac pro motu elliptico, \SourceLeaf{043}praetereaque omnes casus qui forte occurrere possent per methodum aliam infra tradendam absolvere licebit. Postquam $F$ vel $u$ inventa erit, $v$ inde per formulam III, ac dein $r$ vel per II vel per VIII determinabitur; commodius adhuc per formulas VI et VII $v$ et $r$ simul eruentur; e formulis reliquis una alterave pro confirmatione calculi, si lubet, in usum vocari poterit.

\begin{center}
\hypertarget{native-043-article026}{}\pdfbookmark[2]{26.}{article-026}26.
\end{center}
Exemplum. Manentibus $e$ et $b$ ut in exemplo praecedente, sit $t=65{,}41236$: quaeruntur $v$ et $r$. Utendo logarithmis \textsc{Brigg}icis habemus
\begin{center}
\begin{tabular}{@{}ll@{}}
\makebox[18em][l]{$\log t$\hfill .\,.\,.\,.\,.\,.\,.\,.\,.\,.\,.\,.} & $1{,}815\,6598$\\[.3em]
\makebox[18em][l]{$\log\lambda kb^{-\frac{3}{2}}$\hfill .\,.\,.\,.\,.\,.\,.\,.} & $6{,}970\,2758$\\[.3em]
\hline
\makebox[18em][l]{$\log N$\hfill .\,.\,.\,.\,.\,.\,.\,.\,.\,.\,.} & $8{,}785\,9356$,
\end{tabular}\\[.25em]
unde $N=0{,}061\,0851\,4$.
\end{center}
Hinc aequationi $N=\lambda e\operatorname{tang}F-\log\operatorname{tang}(45^0+\frac{1}{2}F)$ satisfieri invenitur per
\[
F=25^0\,24^{\prime}\,27^{\prime\prime}{,}66,
\]
unde fit per formulam III
\begin{center}
\begin{tabular}{@{}ll@{}}
\makebox[18em][l]{$\log\operatorname{tang}\frac{1}{2}F$\hfill .\,.\,.\,.\,.\,.\,.} & $9{,}353\,0120$\\[.3em]
\makebox[18em][l]{$\log\operatorname{tang}\frac{1}{2}\psi$\hfill .\,.\,.\,.\,.\,.\,.} & $9{,}531\,8179$\\[.3em]
\hline
\makebox[18em][l]{$\log\operatorname{tang}\frac{1}{2}v$\hfill .\,.\,.\,.\,.\,.\,.} & $9{,}821\,1941$
\end{tabular}
\end{center}
adeoque $\frac{1}{2}v=33^0\,31^{\prime}\,29^{\prime\prime}{,}89$ atque $v=67^0\,2^{\prime}\,59^{\prime\prime}{,}78$. Hinc porro habetur
\begin{center}
\footnotesize
\begin{tabular}{@{}ll@{}}
$\left.\begin{array}{ll}
\text{C.}\log\cos\frac{1}{2}(v+\psi)\quad .\,.\,. &0{,}213\,7476\\[.3em]
\text{C.}\log\cos\frac{1}{2}(v-\psi)\quad .\,.\,. &0{,}014\,5197
\end{array}\right\}$ & \begin{tabular}{@{}ll@{}}
differentia .\,.\,.\,.\,.\,.\,.\,.\,.\,.\,.\,. & $0{,}199\,2279$\\[.3em]
$\log\operatorname{tang}(45^0+\frac{1}{2}F)$ .\,.\,. & $0{,}199\,2280$
\end{tabular}
\end{tabular}
\end{center}
\begin{center}
\begin{tabular}{@{}ll@{}}
\makebox[18em][l]{$\log\frac{p}{2e}$\hfill .\,.\,.\,.\,.\,.\,.\,.\,.\,.\,.} & $9{,}972\,5868$\\[.3em]
\hline
\makebox[18em][l]{$\log r$\hfill .\,.\,.\,.\,.\,.\,.\,.\,.\,.\,.\,.} & $0{,}200\,8541$
\end{tabular}
\end{center}
